% ============================================================
%  Section 3: Data Modelling for the DWH
% ============================================================
\section{Data Modelling}

% -------- subsection: Logical Modelling Overview --------

\begin{frame}
  \frametitle{Why Data Modelling Matters}

  Data modelling is the \textbf{first and most critical phase} in building a DWH.
  It determines what data is stored, how it is structured, and how easily
  users can analyse it.

  \begin{block}{The Logical Data Model\ldots}
    \begin{itemize}
      \item Is the blueprint of the DWH
      \item Is \textbf{requirements-driven}: what questions does the business need to answer?
      \item Defines which data and in which logical structure is stored in the DWH
      \item Strongly driven by user needs
      \item Is \textbf{critically important for project success}
    \end{itemize}
  \end{block}

  \begin{block}{Two Modelling Phases}
    \begin{enumerate}
      \item \textbf{Logical modelling}: What data and how is it structured?
        (Subject model, dimensional model, ER model)
      \item \textbf{Physical modelling}: How is it technically implemented?
        (Tables, aggregations, partitioning, storage type)
    \end{enumerate}
  \end{block}
\end{frame}

% -------- subsection: Subject Modelling --------

\begin{frame}
  \frametitle{Step 1: Subject Modelling}

  \begin{block}{What is a Subject?}
    A subject is an \textbf{important business area or theme} in the enterprise
    about which a lot of information is collected.
    \begin{itemize}
      \item \textbf{Persons}: Customer, Employee, Supplier
      \item \textbf{Things}: Product, Account, Contract
      \item \textbf{Processes}: Order, Payment, Delivery
    \end{itemize}
  \end{block}

  \begin{block}{The Subject Model\ldots}
    \begin{itemize}
      \item Documents the most important subjects and their relationships
      \item Covers the \textbf{whole enterprise} at a high level
      \item Is easy for business users to understand
      \item Is represented as a graphical diagram (5--15 subjects)
    \end{itemize}
  \end{block}

  \begin{examples}{Tip}
    In operational data models, subjects tend to appear as the \textbf{hub
    of multiple relationships} --- many tables reference them.
    This is a good way to identify them.
  \end{examples}
\end{frame}

% ------------------------------------------------------------
\begin{frame}
  \frametitle{Subjects Vary by Industry}

  \begin{block}{Examples by Sector}
    \begin{tabularx}{\linewidth}{@{}l X@{}}
      \toprule
      \textbf{Sector}    & \textbf{Typical Subjects} \\
      \midrule
      Commercial firm    & Customer, Product, Organisational unit, Sales \\
      Public sector      & Citizen, Department, Employee, Service \\
      University         & Student, Professor, Course, Subject area \\
      Bank               & Customer, Account, Advisor, Product, OU \\
      Insurance          & Customer, Policy, Claim, Agent, Premium \\
      \bottomrule
    \end{tabularx}
  \end{block}

  \begin{alertblock}{Important}
    Subjects must be identified \textbf{before} modelling begins.
    They require analysis of the business requirements and existing data models.
    Use interviews, workshops, and mind maps to elicit them.
  \end{alertblock}
\end{frame}

% ------------------------------------------------------------
\begin{frame}
  \frametitle{Example Subject Model: Bank}

  % PICTURE PLACEHOLDER
  \begin{alertblock}{Picture: Bank Subject Model}
    \textit{Add a subject model graph for a bank with nodes: Customer, Account,
    Product, Advisor, Branch. Show edges (relationships) between them.
    Label the edges (e.g. ``has'', ``belongs to'', ``advises'', ``sells'').
    Keep it simple -- one diagram fitting a slide.}
  \end{alertblock}

  \begin{block}{Subject List Example (Bank)}
    \begin{tabularx}{\linewidth}{@{}l X l@{}}
      \toprule
      \textbf{Subject} & \textbf{Description} & \textbf{Related to} \\
      \midrule
      Customer  & Person/org using bank services & Product, Account \\
      Advisor   & Bank employee serving customers & Branch, Customer \\
      Branch    & Organisational unit of the bank & Advisor, Customer \\
      Account   & Used to manage payment transfers & Customer, Product \\
      Product   & Services sold to customers & Customer, Account \\
      \bottomrule
    \end{tabularx}
  \end{block}
\end{frame}

% -------- subsection: Dimensional Modelling --------

\begin{frame}
  \frametitle{Step 2: Dimensional Modelling}

  \begin{block}{The Dimensional Data Model\ldots}
    \begin{itemize}
      \item Provides a \textbf{detailed logical view} of the data
      \item Divides data into \textbf{Facts} and \textbf{Dimensions}
      \item Is \textbf{requirements-oriented}: what analyses are needed?
      \item Is represented like an Entity-Relationship diagram
      \item The key difference from classical ER: the \textbf{degree of normalisation}
        is deliberately lower (denormalisation for query performance)
    \end{itemize}
  \end{block}

  \begin{block}{Core Concepts}
    \begin{tabularx}{\linewidth}{@{}l X@{}}
      \textbf{Facts}      & The measurable business metrics to be analysed \\
      \textbf{Dimensions} & The perspectives from which facts are analysed \\
      \textbf{Grain}      & The lowest level of detail stored (most important!) \\
    \end{tabularx}
  \end{block}
\end{frame}

% ------------------------------------------------------------
\begin{frame}
  \frametitle{Facts}

  \begin{block}{Definition}
    Facts contain the \textbf{business metrics (KPIs)} to be analysed.
  \end{block}

  \begin{block}{Sources and Types}
    \begin{itemize}
      \item Either taken directly from operational systems:
        quantities, prices, measurements
      \item Or derived from them: revenue, profit margin, averages
      \item Stored at the \textbf{lowest level of detail} (finest grain)
        in the DWH
    \end{itemize}
  \end{block}

  \begin{block}{Additive Types of Facts}
    \begin{tabularx}{\linewidth}{@{}l X l@{}}
      \toprule
      \textbf{Type}     & \textbf{Description}                             & \textbf{Example} \\
      \midrule
      Additive          & SUM over all dimensions                          & Revenue, orders \\
      Semi-additive     & SUM over some dimensions only                    & Account balance \\
      Non-additive      & Cannot be meaningfully summed                    & Ratios, prices \\
      \bottomrule
    \end{tabularx}
  \end{block}
\end{frame}

% ------------------------------------------------------------
\begin{frame}
  \frametitle{Dimensions}

  \begin{block}{Definition}
    Dimensions are the \textbf{logical perspectives} under which facts
    are analysed.
  \end{block}

  \begin{itemize}
    \item Logical groupings of attributes: e.g.\ Year, Month, Day $\to$ \textbf{Time}
    \item Can contain additional descriptive attributes (e.g.\ customer first name)
    \item Each dimension gets its own table
    \item Dimensions often correspond to the \textbf{subjects} from the subject model
    \item Connected to the fact table via \textbf{dimension keys}
  \end{itemize}

  \begin{examples}{Typical Dimensions}
    \textbf{Time} (Year/Quarter/Month/Day),
    \textbf{Product} (Category/Group/Individual Product),
    \textbf{Customer} (Segment/Region/Individual),
    \textbf{Geography} (Continent/Country/Region/City),
    \textbf{Organisation} (Company/Division/Department)
  \end{examples}
\end{frame}

% ------------------------------------------------------------
\begin{frame}
  \frametitle{Dimension Hierarchies}

  Dimensions are organised in \textbf{hierarchies} --- levels of detail from
  coarse to fine. Hierarchies are the basis for drill-down and roll-up.

  \begin{columns}[T]
    \column{0.48\textwidth}
    \begin{block}{Product Hierarchy}
      \begin{itemize}
        \item Total
        \item Industry
        \item Product Group
        \item Product Family
        \item Product (lowest level)
      \end{itemize}
    \end{block}

    \column{0.48\textwidth}
    \begin{block}{Time Dimension: Multiple Hierarchies}
      \begin{itemize}
        \item \textit{Hierarchy 1:} Year $\to$ Quarter $\to$ Month $\to$ Day
        \item \textit{Hierarchy 2:} Year $\to$ 4-4-5 Month $\to$ Week $\to$ Day
      \end{itemize}
      One dimension can contain \textbf{multiple} hierarchies.
    \end{block}
  \end{columns}

  % PICTURE PLACEHOLDER
  \begin{alertblock}{Picture: Product Dimension Hierarchy}
    \textit{Add a tree/pyramid diagram for the Product dimension:
    Gesamt at top, then Branche, Produktgruppe, Produktfamilie, Produkt at bottom.
    Arrows pointing downward. Label it ``Drill-Down'' direction.}
  \end{alertblock}
\end{frame}

% ------------------------------------------------------------
\begin{frame}
  \frametitle{The Multidimensional Data Model}

  The dimensional model conceptually organises data as a
  \textbf{multidimensional space} (hypercube).

  % PICTURE PLACEHOLDER
  \begin{alertblock}{Picture: 3D Multidimensional Cube}
    \textit{Add a 3D cube diagram with three axes:
    Dimension 1 = Organisation (Branch/Region/Advisor),
    Dimension 2 = Product Line (Group/Individual),
    Dimension 3 = Time (Year/Month/Day).
    Highlight one cell with a value, e.g. ``Org=Branch02, Product=03, Time=1999''.}
  \end{alertblock}

  \begin{block}{Key Concepts}
    \begin{itemize}
      \item Each \textbf{cell} in the cube holds an aggregated measure
      \item \textbf{Sparse cubes}: many cells have no data (normal in practice)
      \item \textbf{Dense cubes}: most cells contain values
      \item Real cubes can have many more than 3 dimensions
    \end{itemize}
  \end{block}
\end{frame}

% -------- subsection: Star Schema --------

\begin{frame}
  \frametitle{The Star Schema}

  \begin{block}{Structure}
    \begin{itemize}
      \item One central \textbf{fact table} connected to multiple
        \textbf{dimension tables} via foreign keys
      \item Dimension tables are \textbf{denormalised} (flat, wide tables)
      \item Simple to understand and query
      \item Optimised for read performance
      \item Named for its star-like visual shape
    \end{itemize}
  \end{block}

  \begin{examples}{Example: Wine Sales}
    Fact table: \texttt{FACT\_SALES} with columns Product-Key,
    Region-Key, Taste-Key, Time-Key, Revenue.\\
    Dimension tables: \texttt{DIM\_PRODUCT} (Rosé/White/Sparkling),
    \texttt{DIM\_REGION} (Brandenburg/Hessen/\ldots),
    \texttt{DIM\_TASTE} (dry/semi-dry/mild),
    \texttt{DIM\_TIME} (Jan/Feb/Mar/\ldots)
  \end{examples}
\end{frame}

% ------------------------------------------------------------
\begin{frame}
  \frametitle{Star Schema: Diagram}

  % PICTURE PLACEHOLDER
  \begin{alertblock}{Picture: Star Schema Diagram (Wine Sales)}
    \textit{Add the wine sales star schema diagram:
    Central fact table FACT\_SALES with foreign keys to four dimension tables.
    DIM\_PRODUCT on the left (Rotwein, Weißwein, Sekt),
    DIM\_REGION on the right (Brandenburg, Hessen, Niedersachsen),
    DIM\_GESCHMACK below-left (trocken, halbtrocken, mild),
    DIM\_ZEITRAUM below-right (Januar, Februar, März).
    Connect all four to the central fact table.}
  \end{alertblock}

  \begin{block}{Fact Table Columns}
    \begin{itemize}
      \item \textbf{Dimension keys} (FK): link to each dimension table
      \item \textbf{Degenerate dimensions}: attributes with no dimension table
        (e.g.\ invoice number)
      \item \textbf{Measures}: the actual facts (revenue, quantity, \ldots)
      \item \textbf{Derived values}: calculated metrics
    \end{itemize}
  \end{block}
\end{frame}

% -------- subsection: Snowflake Schema --------

\begin{frame}
  \frametitle{The Snowflake Schema}

  \begin{block}{Concept}
    The snowflake schema is a \textbf{normalised extension} of the star schema.
    Large or redundant dimension tables are split into multiple related tables.
  \end{block}

  \begin{examples}{Example: Region Dimension}
    In a star schema, one flat \texttt{DIM\_REGION} table contains
    (City, PostalCode, FederalState, Country).

    In a snowflake schema, this splits into:
    \texttt{DIM\_REGION} $\to$ \texttt{DIM\_POSTAL\_AREA} $\to$ \texttt{DIM\_CITY}
  \end{examples}

  % PICTURE PLACEHOLDER
  \begin{alertblock}{Picture: Snowflake Schema Diagram}
    \textit{Add the snowflake extension of the wine sales star schema:
    DIM\_REGION now has a FK to DIM\_POSTAL\_AREA, which has a FK to DIM\_CITY.
    DIM\_ZEITRAUM splits into Month-ID and Week-ID sub-tables.
    Show the ``snowflake'' branching effect.}
  \end{alertblock}
\end{frame}

% ------------------------------------------------------------
\begin{frame}
  \frametitle{Star vs.\ Snowflake Schema}

  \begin{tabularx}{\linewidth}{@{}l X X@{}}
    \toprule
    \textbf{Property}    & \textbf{Star Schema}        & \textbf{Snowflake Schema} \\
    \midrule
    Normalisation        & Denormalised                & Partially normalised \\
    Query simplicity     & Simple (fewer joins)        & More complex (more joins) \\
    Query performance    & Faster                      & Slightly slower \\
    Storage redundancy   & Higher                      & Lower \\
    Maintenance          & Easier                      & More complex \\
    Typical use          & Smaller dimensions          & Very large dimensions \\
    \bottomrule
  \end{tabularx}

  \vspace{0.5em}
  \begin{block}{In practice\ldots}
    The star schema is the \textbf{most common choice} for data marts.
    The snowflake schema is used selectively to reduce storage for
    specific very large dimension tables.
    \textbf{Most real-world DWH schemas are hybrid.}
  \end{block}
\end{frame}

% -------- subsection: Date Dimension --------

\begin{frame}
  \frametitle{The Date Dimension}

  The time/date dimension is present in \textbf{almost every DWH}.
  It is pre-populated for all dates in the analysis range (e.g.\ 2000--2040).

  \begin{examples}{Typical Date Dimension Table}
    \scriptsize
    \begin{tabularx}{\linewidth}{@{}l l l l l l l l@{}}
      \toprule
      \textbf{Date}    & \textbf{Weekday} & \textbf{Month} & \textbf{Cal.Q} & \textbf{Fin.Q} & \textbf{Week} & \textbf{4-4-5} & \textbf{Holiday} \\
      \midrule
      25.12.1997 & Thu & 12 & 4 & 1 & 52 & 4 & Yes \\
      26.12.1997 & Fri & 12 & 4 & 1 & 52 & 4 & Yes \\
      31.12.1997 & Wed & 12 & 4 & 1 & 53 & 4 & No \\
      01.01.1998 & Thu & 01 & 1 & 1 & 01 & 1 & Yes \\
      \bottomrule
    \end{tabularx}
    \normalsize
  \end{examples}

  \begin{alertblock}{Important Note}
    Time is \textbf{not} an isolated dimension. Every other dimension
    can also change over time (available products, customer lists, sales territories).
    This is why \textbf{Slowly Changing Dimensions} are crucial.
  \end{alertblock}
\end{frame}

% -------- subsection: SCD --------

\begin{frame}
  \frametitle{Slowly Changing Dimensions (SCD)}

  \begin{block}{The Problem}
    Dimension data changes over time. A customer moves cities.
    A product is re-categorised. A sales region is restructured.
    How do we handle this in a non-volatile DWH?
  \end{block}

  \begin{block}{SCD Type 1 --- Overwrite}
    Simply \textbf{update} the existing row. No history is retained.\\
    \textit{Use when historical accuracy for this attribute does not matter.}\\
    Example: Fixing a data entry error.
  \end{block}

  \begin{block}{SCD Type 2 --- Add a New Row (most common)}
    Insert a \textbf{new row} for the new value. Old rows remain unchanged
    with a validity period (start/end date or a \texttt{is\_current} flag).\\
    \textit{Use when you need to track the history of changes.}\\
    Example: Customer changes address.
  \end{block}

  \begin{block}{SCD Type 3 --- Add a New Column}
    Add a column for the \textbf{previous value}.\\
    \textit{Only supports one level of history.}
  \end{block}
\end{frame}

% -------- subsection: Dimension Keys --------

\begin{frame}
  \frametitle{Dimension Keys: Natural vs.\ Surrogate Keys}

  \begin{block}{Natural Key}
    The business identifier from the source system
    (e.g.\ customer number, product code, employee ID).
  \end{block}

  \begin{block}{Surrogate Key}
    A \textbf{meaningless integer} generated by the DWH (e.g.\ 1, 2, 3, \ldots).
    It is the primary key of the dimension table inside the DWH.
  \end{block}

  \begin{block}{Why Use Surrogate Keys?}
    \begin{itemize}
      \item Natural keys can change in source systems --- surrogate keys are stable
      \item Essential for \textbf{SCD Type 2}: multiple rows for the same business entity
        get different surrogate keys
      \item Smaller (integer) = faster joins
      \item Decouples the DWH from source system changes
      \item Allows integration of data from \textbf{multiple source systems}
        (different natural keys for the same entity)
    \end{itemize}
  \end{block}
\end{frame}

% -------- subsection: ER Modelling --------

\begin{frame}
  \frametitle{Entity-Relationship Modelling in the DWH Context}

  \begin{block}{When ER is Used}
    ER modelling (Peter Chen, 1976) is the traditional approach for
    relational data modelling. In the DWH context, it is used:
    \begin{itemize}
      \item For modelling the \textbf{integration/staging layer} (3NF)
      \item Favoured by the Inmon approach
      \item Particularly popular with database and IT specialists
    \end{itemize}
  \end{block}

  \begin{block}{Core ER Concepts}
    \begin{tabularx}{\linewidth}{@{}l X@{}}
      \textbf{Entity}      & A distinguishable object (Customer, Product) \\
      \textbf{Attribute}   & A property of an entity (Name, Price) \\
      \textbf{Relationship}& A connection between two entities \\
      \textbf{Cardinality} & How many instances participate: 1:1, 1:n, n:m \\
      \textbf{Optionality} & Whether participation is mandatory or optional \\
    \end{tabularx}
  \end{block}
\end{frame}

% ------------------------------------------------------------
\begin{frame}
  \frametitle{Normalisation in the DWH}

  \begin{block}{Why Normalise?}
    \begin{itemize}
      \item Avoids redundancy and inconsistency
      \item Prevents update/insert/delete anomalies
      \item Results in a clearer, more maintainable data model
      \item Concept introduced by E.F. Codd (1972)
    \end{itemize}
  \end{block}

  \begin{block}{Normal Forms Recap}
    \begin{tabularx}{\linewidth}{@{}l X@{}}
      \toprule
      \textbf{Form} & \textbf{Rule} \\
      \midrule
      1NF & All attribute values are atomic (no repeating groups) \\
      2NF & 1NF + every non-key attribute fully depends on the whole PK \\
      3NF & 2NF + no non-key attribute depends on another non-key attribute \\
      BCNF & Stronger form of 3NF \\
      \bottomrule
    \end{tabularx}
  \end{block}

  \begin{alertblock}{DWH vs.\ OLTP}
    The core DWH/staging layer uses \textbf{3NF} for clean data storage.
    Data marts use \textbf{deliberately denormalised} star/snowflake schemas
    for analytical query performance.
  \end{alertblock}
\end{frame}
